\setcounter{framenumber}{0}
%29

\begin{frame}
  \titlepage
\end{frame}

\begin{frame}{Learning Goals}
By the end of this lecture, students should be able to:
\begin{itemize}
    \item apply the Cauchy--Schwarz inequality;
    \item use Jensen's inequality with convex and concave functions;
    \item apply Markov's inequality;
    \item apply Chebyshev's inequality;
    \item apply Chernoff bounds;
    \item interpret inequalities probabilistically.
\end{itemize}
\end{frame}

\lecturepart{Why Inequalities Matter}

\begin{frame}{What If Exact Calculation Is Impossible?}
In many real problems:
\begin{itemize}
    \item exact probabilities are difficult;
    \item integrals may be impossible;
    \item distributions may be unknown;
    \item computations may be too expensive.
\end{itemize}

\vspace{0.25cm}

Three major strategies:
\begin{enumerate}
    \item simulation;
    \item inequalities (bounds);
    \item asymptotic approximations.
\end{enumerate}

\vspace{0.25cm}

\begin{keybox}{Today's focus}
Probability inequalities
\end{keybox}
\end{frame}

\begin{frame}{What Is a Bound?}
Suppose we cannot compute
\[
\Prob(A)
\]
exactly.

\vspace{0.25cm}

An inequality may tell us:
\[
0.01
\le
\Prob(A)
\le
0.03.
\]

\vspace{0.25cm}

Even without the exact answer, we gain useful information.

\end{frame}

\lecturepart{Cauchy--Schwarz Inequality}

\begin{frame}{Cauchy--Schwarz Inequality}
\begin{keybox}{Theorem}
For random variables $X,Y$ with finite second moments,
\[
{
|\mathbb{E}(XY)|
\le
\sqrt{
\mathbb{E}(X^2)\mathbb{E}(Y^2)
}.
}
\]
\end{keybox}

\vspace{0.25cm}

This is one of the most applicable inequalities in mathematics (linear algebra, statistics, functional inequalities, geometry, etc.).
\end{frame}

\begin{frame}{Idea Behind the Proof}
For every real number $t$,
\[
(Y-tX)^2\ge0.
\]

Taking expectations:
\[
\mathbb{E}(Y-tX)^2\ge0.
\]

Expanding:
\[
\mathbb{E}(Y^2)
-2t\mathbb{E}(XY)
+t^2\mathbb{E}(X^2)
\ge0.
\]

\vspace{0.25cm}

This quadratic must always be nonnegative.
\end{frame}

\begin{frame}{Proof of Cauchy--Schwarz}
Since
\[
\mathbb{E}(Y^2)
-2t\mathbb{E}(XY)
+t^2\mathbb{E}(X^2)
\ge0
\]
for all $t$,
its discriminant must satisfy:
\[
(-2\mathbb{E}(XY))^2
-
4\mathbb{E}(X^2)\mathbb{E}(Y^2)
\le0.
\]

Thus,
\[
\mathbb{E}(XY)^2
\le
\mathbb{E}(X^2)\mathbb{E}(Y^2).
\]

Taking square roots gives:
\[
{
|\mathbb{E}(XY)|
\le
\sqrt{
\mathbb{E}(X^2)\mathbb{E}(Y^2)
}.
}
\]
\end{frame}

\begin{frame}{Application of Cauchy-Schwarz: $-1\leq Corr(X,Y)\leq 1$}
Apply Cauchy--Schwarz to
\[
X-\mathbb{E}(X)
\]
and
\[
Y-\mathbb{E}(Y).
\]

Then:
\[
|\operatorname{Cov}(X,Y)|
\le
\sqrt{
\operatorname{Var}(X)\operatorname{Var}(Y)
}.
\]

Dividing both sides:
\[
{
|\operatorname{Corr}(X,Y)|
\le1.
}
\]
\end{frame}

\begin{frame}{Application of Cauchy-Schwarz: $Var(X)\geq 0$}
Apply Cauchy--Schwarz with
\[
Y=1.
\]

Then:
\[
|\mathbb{E}(X)|^2
\le
\mathbb{E}(X^2).
\]

Therefore,
\[
\boxed{
\operatorname{Var}(X)
=
\mathbb{E}(X^2)-(\mathbb{E}(X))^2
\ge0.
}
\]
\end{frame}


\lecturepart{Jensen's Inequality}
\begin{frame}{Convex Functions}
\begin{keybox}{Definition}
A function \(g:I\to\mathbb{R}\) is called convex if
\[
g(\lambda x+(1-\lambda)y)
\le
\lambda g(x)+(1-\lambda)g(y)
\]
for all \(x,y\in I\) and all
\[
0\le \lambda \le 1.
\]
\end{keybox}


Interpretation: $\lambda x+(1-\lambda)y$ is a weighted average of \(x\) and \(y\), and convexity says $g(\text{average})
\le
\text{average of } g.$

\textbf{Geometrically}: the graph of a convex function always lies below its secant lines.

\begin{keybox}{Second Derivative Test}
If \(g\) is twice differentiable, then:
\[
g''(x)\ge0
\quad \text{for all } x
\]
implies that \(g\) is convex.
\end{keybox}
 

Examples of convex functions:
\[
g(x)=x^2,
\qquad
g(x)=e^x.
\]
\end{frame}

\begin{frame}{Jensen's Inequality}
\begin{keybox}{Theorem}
If $g$ is convex, then
\[
{
\mathbb{E}(g(X))
\ge
g(\mathbb{E}(X)).
}
\]
\end{keybox}

If $g$ is concave, the inequality reverses:
\[
{
\mathbb{E}(g(X))
\le
g(\mathbb{E}(X)).
}
\]

The theorem is proved using induction (not covered - available in Wikipedia).
\end{frame}

\begin{frame}{Geometric Interpretation of Jensen}
For a convex function:
\begin{itemize}
    \item tangent lines lie below the graph;
    \item averaging inputs produces smaller values than averaging outputs.
\end{itemize}

\vspace{0.25cm}

\begin{keybox}{Mnemonic}
\[
g(\text{average})
\le
\text{average of }g.
\]
for convex functions.
\end{keybox}
\end{frame}

\begin{frame}{Example (Jensen's Inequality): $Var(X)\geq 0 $ (again)}
Take
\[
g(x)=x^2.
\]

Since
\[
g''(x)=2>0,
\]
the function is convex.

\vspace{0.25cm}

Jensen gives:
\[
\mathbb{E}(X^2)
\ge
(\mathbb{E}(X))^2.
\]

Thus:
\[
\operatorname{Var}(X)\ge0.
\]
\end{frame}

\begin{frame}{Example (Jensen's Inequality): Logarithms}
Take
\[
g(x)=\log x.
\]

Since
\[
g''(x)=-\frac1{x^2}<0,
\]
the logarithm is concave.

\vspace{0.25cm}

Thus:
\[
{
\mathbb{E}(\log X)
\le
\log(\mathbb{E}(X)).
}
\]

\vspace{0.25cm}

This inequality is fundamental in statistics and information theory.
\end{frame}

\begin{frame}{Example (Jensen's Inequality): Sample Standard Deviation Is Biased}
Recall:
\[
\mathbb{E} [S_n^2] = \sigma^2, 
\]
i.e., $S^2$ is an unbiased estimator for $\sigma^2$.

Now, since
\[
g(x)=\sqrt{x}
\]
is concave,
Jensen's inequality gives:
\[
\mathbb{E}(S_n)
=
\mathbb{E}(\sqrt{S_n^2})
\le
\sqrt{\mathbb{E}(S_n^2)}
=
\sigma.
\]

\vspace{0.25cm}

\begin{keybox}{Conclusion}
Sample standard deviation tends to underestimate the true standard deviation.
\end{keybox}
\end{frame}

\lecturepart{Tail Bounds}

\begin{frame}{Tail Probabilities}
We often want bounds on:
\[
\Prob(X\ge a),
\]
or
\[
\Prob(|X-\mu|\ge a).
\]

\vspace{0.25cm}

These are called tail probabilities.

\vspace{0.25cm}

\begin{keybox}{Applications}
\begin{itemize}
    \item rare events;
    \item risk analysis;
    \item statistics;
    \item machine learning;
    \item concentration inequalities.
\end{itemize}
\end{keybox}
\end{frame}

\begin{frame}{Markov's Inequality}
\begin{keybox}{Theorem (Markov's Inequality)}
If $X\ge0$ and $a>0,$ then
\[
{
\Prob(X\ge a)
\le
\frac{\mathbb{E}(X)}{a}.
}
\]
\end{keybox}

\textbf{Proof}

First, note that $X
\ge
a\,I(X\ge a).$

Taking expectations from both sides,
\[
\mathbb{E}(X)
\ge
a\mathbb{E}(I(X\ge a)).
\]

Using the fundamental bridge:
\[
\mathbb{E}(I(X\ge a))
=
\Prob(X\ge a).
\]

Therefore:
\[
{
\Prob(X\ge a)
\le
\frac{\mathbb{E}(X)}{a}.
}
\]

\end{frame}



\begin{frame}{Chebyshev's Inequality}
\begin{keybox}{Theorem (Chebyshev's Inequality)}
If $X$ has mean $\mu$ and variance $\sigma^2$, then for every $a>0,$
\[
{
\Prob(|X-\mu|\ge a)
\le
\frac{\sigma^2}{a^2}.
}
\]
\end{keybox}
\textbf{Proof}

Apply Markov to $(X-\mu)^2\geq 0.$ Then,
\[
\Prob((X-\mu)^2\ge a^2)
\le
\frac{
\mathbb{E}(X-\mu)^2
}{
a^2 
}=\frac{\sigma^2}{a^2}.
\]
\end{frame}



\begin{frame}{Standardized Form of Chebyshev}
Let
\[
a=c\sigma.
\]

Then:
\[
\boxed{
\Prob(|X-\mu|\ge c\sigma)
\le
\frac1{c^2}.
}
\]

\vspace{0.25cm}

Examples:
\[
\Prob(|X-\mu|\ge2\sigma)\le\frac14,
\]

\[
\Prob(|X-\mu|\ge3\sigma)\le\frac19.
\]
\end{frame}

\begin{frame}{Chernoff Bounds}
Markov becomes stronger if we apply it to:
\[
e^{tX},
\qquad t>0.
\]

\vspace{0.25cm}

\begin{keybox}{Chernoff Bound}
For every
\[
t>0,
\]
\[
\boxed{
\Prob(X\ge a)
\le
\frac{
\mathbb{E}(e^{tX})
}{
e^{ta}
}.
}
\]
\end{keybox}

\vspace{0.25cm}

The numerator is the MGF of $X$.
\end{frame}
 

\begin{frame}{Comparing Bounds for Normal Tails}
Suppose
\[
Z\sim N(0,1).
\]

True probability:
\[
\Prob(|Z|>3)\approx0.003.
\]

Bounds:
\[
\text{Markov: }0.27,
\]

\[
\text{Chebyshev: }0.11,
\]

\[
\text{Chernoff: }0.022.
\]

\vspace{0.25cm}

\begin{keybox}{Observation}
Chernoff bounds are often dramatically better.
\end{keybox}
\end{frame}

\begin{frame}{Summary}
\begin{itemize}
    \item Cauchy--Schwarz bounds expectations.

    \item Jensen compares
    \[
    \mathbb{E}(g(X))
    \]
    and
    \[
    g(\mathbb{E}(X)).
    \]

    \item Markov bounds nonnegative tails.

    \item Chebyshev controls deviations from the mean.

    \item Chernoff bounds are often much sharper.

    \item Inequalities are fundamental tools throughout probability and statistics.
\end{itemize}

\vspace{0.25cm}

\begin{keybox}{Next lecture}
Law of Large Numbers and Central Limit Theorem.
\end{keybox}
\end{frame}
